% GENERATED FILE. Edit canonical lessons, then run npm run generate:books.
% canonical-source-hash: fnv1a64:cb18fcb5962cb622
% canonical-lessons: IT-C411-cofano, IT-C411-bagagliaio, IT-C411-frizione, IT-C411-fanale, IT-C411-svitare

\chapter{Sotto il cofano}
\label{ch:it-sotto-il-cofano}
\hlchaptermodality{411}

\begin{chapteropening}
\textbf{By the end of this chapter:} \emph{I can talk about parts of a car: the bonnet, the boot, the clutch, the lights, and unscrewing a wheel.}

Complete the last lesson of chapter 411: I can talk about parts of a car: the bonnet, the boot, the clutch, the lights, and unscrewing a wheel.
\end{chapteropening}

\section[il cofano]{il cofano — the bonnet (of a car); the chest}

\label{lesson:IT-C411-cofano}



\begin{quote}
Before the new one: say the Italian for to tow, then the Italian for the garage.
\end{quote}

\subsection*{You'll want to know: il cofano}
\textbf{il cofano} — \textquotedblleft{}the bonnet (of a car); the chest\textquotedblright{}.

\begin{quote}
Il meccanico apre il cofano.
\end{quote}

\begin{quote}
\emph{The mechanic opens the bonnet.}
\end{quote}

In American English, the hood. Its old sense is a chest or coffer.

\begin{cousinweb}
\emph{cofano} is inherited from Latin \emph{cophinus} 'basket' (EtymDB), from Greek \emph{kóphinos}; it first meant a chest or coffer.
\end{cousinweb}

\begin{grammarlens}[title={What you've built}]
One more word for telling, explaining and describing at B1.
\end{grammarlens}

\subsection*{Your turn}
\begin{itemize}
\raggedright
  \item \emph{Say it:} \emph{il cofano}
  \item \emph{Say it:} \emph{il cofano}, once more
  \item \emph{Say it:} a sentence of your own with \emph{il cofano}
  \item \emph{Recall:} say the Italian for to tow, then the Italian for the garage, then say \emph{il cofano} again
\end{itemize}

\begin{tcolorbox}[breakable,colback=teal!4,colframe=teal!35!black,title={Before you move on}]
What does il cofano mean? (\textquotedblleft{}The bonnet (of a car); the chest\textquotedblright{}.) Say it once more.
\end{tcolorbox}
\section[il bagagliaio]{il bagagliaio — the boot (of a car)}

\label{lesson:IT-C411-bagagliaio}



\begin{quote}
Before the new one: say the Italian for the garage, then the Italian for the bonnet.
\end{quote}

\subsection*{You'll want to know: il bagagliaio}
\textbf{il bagagliaio} — \textquotedblleft{}the boot (of a car)\textquotedblright{}.

\begin{quote}
Metti le valigie nel bagagliaio.
\end{quote}

\begin{quote}
\emph{Put the suitcases in the boot.}
\end{quote}

In American English, the trunk. \emph{Il bagaglio} is luggage.

\begin{cousinweb}
\emph{bagagliaio} is built on \emph{bagaglio} 'luggage', which came from French \emph{bagage}, built on Old French \emph{bague} 'bundle'; where that word comes from is not settled.
\end{cousinweb}

\begin{grammarlens}[title={What you've built}]
One more word for telling, explaining and describing at B1.
\end{grammarlens}

\subsection*{Your turn}
\begin{itemize}
\raggedright
  \item \emph{Say it:} \emph{il bagagliaio}
  \item \emph{Say it:} \emph{il bagagliaio}, once more
  \item \emph{Say it:} a sentence of your own with \emph{il bagagliaio}
  \item \emph{Recall:} say the Italian for the garage, then the Italian for the bonnet, then say \emph{il bagagliaio} again
\end{itemize}

\begin{tcolorbox}[breakable,colback=teal!4,colframe=teal!35!black,title={Before you move on}]
What does il bagagliaio mean? (\textquotedblleft{}The boot (of a car)\textquotedblright{}.) Say it once more.
\end{tcolorbox}
\section[la frizione]{la frizione — the clutch; the rub, the rubdown}

\label{lesson:IT-C411-frizione}



\begin{quote}
Before the new one: say the Italian for the bonnet, then the Italian for the boot.
\end{quote}

\subsection*{You'll want to know: la frizione}
\textbf{la frizione} — \textquotedblleft{}the clutch; the rub, the rubdown\textquotedblright{}.

\begin{quote}
Premi la frizione e cambia marcia.
\end{quote}

\begin{quote}
\emph{Press the clutch and change gear.}
\end{quote}

\emph{Frizione} also means friction, and a rubdown with lotion.

\begin{cousinweb}
\emph{frizione} is a learned borrowing from Latin \emph{frictiō} 'rubbing' (EtymDB), from \emph{fricāre} 'to rub': the clutch plates rub together.
\end{cousinweb}

\begin{grammarlens}[title={What you've built}]
One more word for telling, explaining and describing at B1.
\end{grammarlens}

\subsection*{Your turn}
\begin{itemize}
\raggedright
  \item \emph{Say it:} \emph{la frizione}
  \item \emph{Say it:} \emph{la frizione}, once more
  \item \emph{Say it:} a sentence of your own with \emph{la frizione}
  \item \emph{Recall:} say the Italian for the bonnet, then the Italian for the boot, then say \emph{la frizione} again
\end{itemize}

\begin{tcolorbox}[breakable,colback=teal!4,colframe=teal!35!black,title={Before you move on}]
What does la frizione mean? (\textquotedblleft{}The clutch; the rub, the rubdown\textquotedblright{}.) Say it once more.
\end{tcolorbox}
\section[il fanale]{il fanale — the light (of a car), the lamp}

\label{lesson:IT-C411-fanale}



\begin{quote}
Before the new one: say the Italian for the boot, then the Italian for the clutch.
\end{quote}

\subsection*{You'll want to know: il fanale}
\textbf{il fanale} — \textquotedblleft{}the light (of a car), the lamp\textquotedblright{}.

\begin{quote}
Un fanale della moto è rotto.
\end{quote}

\begin{quote}
\emph{One of the motorbike lights is broken.}
\end{quote}

\emph{Il fanale posteriore} — the rear light.

\begin{cousinweb}
\emph{fanale} is usually traced to Greek \emph{phanós} 'lantern, torch', borrowed through Byzantine Greek, the Greek of the eastern empire.
\end{cousinweb}

\begin{grammarlens}[title={What you've built}]
One more word for telling, explaining and describing at B1.
\end{grammarlens}

\subsection*{Your turn}
\begin{itemize}
\raggedright
  \item \emph{Say it:} \emph{il fanale}
  \item \emph{Say it:} \emph{il fanale}, once more
  \item \emph{Say it:} a sentence of your own with \emph{il fanale}
  \item \emph{Recall:} say the Italian for the boot, then the Italian for the clutch, then say \emph{il fanale} again
\end{itemize}

\begin{tcolorbox}[breakable,colback=teal!4,colframe=teal!35!black,title={Before you move on}]
What does il fanale mean? (\textquotedblleft{}The light (of a car), the lamp\textquotedblright{}.) Say it once more.
\end{tcolorbox}
\section[svitare]{svitare — to unscrew}

\label{lesson:IT-C411-svitare}



\begin{quote}
Before the new one: say the Italian for the clutch, then the Italian for the light.
\end{quote}

\subsection*{You'll want to know: svitare}
\textbf{svitare} — \textquotedblleft{}to unscrew\textquotedblright{}.

\begin{quote}
Svita la ruota con la chiave.
\end{quote}

\begin{quote}
\emph{Unscrew the wheel with the wrench.}
\end{quote}

The opposite is \emph{avvitare}, to screw in.

\begin{cousinweb}
\emph{svitare} is built from \emph{s-} and \emph{vite} 'screw', which is inherited from Latin \emph{vītis} 'vine' (EtymDB): the thread of a screw winds like a vine.
\end{cousinweb}

\begin{grammarlens}[title={What you've built}]
One more word for telling, explaining and describing at B1.
\end{grammarlens}

\subsection*{Your turn}
\begin{itemize}
\raggedright
  \item \emph{Say it:} \emph{svitare}
  \item \emph{Say it:} \emph{svitare}, once more
  \item \emph{Say it:} a sentence of your own with \emph{svitare}
  \item \emph{Recall:} say the Italian for the clutch, then the Italian for the light, then say \emph{svitare} again
\end{itemize}

\begin{tcolorbox}[breakable,colback=teal!4,colframe=teal!35!black,title={Before you move on}]
What does svitare mean? (\textquotedblleft{}To unscrew\textquotedblright{}.) Say it once more.
\end{tcolorbox}
